%% kapa: exam=1 round=23 year=2012 period=2 form=A subject=appraisal_law subject_ko=부동산관계법규 seq=10 number=50
%% source: Q-Net 감정평가사 기출문제 (cst00309) artlSeq=5052580 file=raw/1008/5052580/2012년 제23회 감정평가사 제1차시험 기출문제(2교시 A형).pdf
%% conversion: kordoc 4.24.51 -> tools/gen1.py (auto); review=proofread-sonnet; flags=none
\begin{question}{50}
국토의 계획 및 이용에 관한 법령상 토지거래허가에 관한 설명으로 옳지 않은 것은?
\begin{choices}
  \item 허가신청 면적이 그 토지의 이용목적으로 보아 적합하지 않다고 인정되는 경우에는 허가를 받을 수 없다.
  \item 기존의 농업인이 아닌 자가 농업용 토지를 취득하고자 하는 경우 그가 거주하는 시ㆍ군에 소재하는 토지와 거주지로부터 20 km 이내에 소재하는 토지에 한해서 허가를 받을 수 있다.
  \item 토지거래계약을 체결하려는 자의 토지이용목적이 도시계획에 맞지 아니한 경우에는 허가를 받을 수 없다.
  \item 법률에 따라 토지를 수용할 수 있는 사업을 시행하는 자가 그 사업을 시행하기 위하여 필요한 것인 경우 허가를 받을 수 있다.
  \item 허가구역을 포함한 지역의 주민을 위한 편익시설로서 관할 시장ㆍ군수 또는 구청장이 확인한 시설의 설치에 이용하려는 경우 허가를 받을 수 있다.
\end{choices}
\end{question}
\answer{2}
